\chapter{Решение ветви: от непрерывной нити к спектральному примитиву перехода}
\label{ch:v6-single-thread-scale-hierarchy-branch-decision-gate}

\section{Причина архитектурного решения}

Предыдущий гейт установил, что текущий родитель не содержит
фундаментального радиуса нити, положительного reach, исключённого объёма
или бесконечного барьера пересоединения. Оставалась последняя допустимая
лазейка: возможно, cutoff, конечный носитель или внутренний радиус уже
задают скрытую иерархию $a\ll\ell$, превращающую раннюю верёвку в
дискретную ткань без нового параметра.

Полный ledger масштабов даёт отрицательный ответ:
\begin{enumerate}
 \item топология и нормированная конечная геометрия фиксируют только
 безразмерные отношения;
 \item спектральное действие сохраняет независимые $f_2,f_0,\Lambda$ и
 внутреннюю длину;
 \item профиль бозонного дефекта сохраняет независимые единицы $L_0,E_0$;
 \item формальное планковское сопоставление Тома IV нарушает собственные
 условия локального heat-kernel разложения;
 \item hard-core правило не найдено.
\end{enumerate}

Следовательно, буквальную непрерывную неразрывную нить нельзя продолжать
как следствие текущего родителя. Она сохраняется только как историческая
и крупномасштабная метафора.

\section{Уже существующий дискретный кандидат}

Отказ от верёвки не требует вводить новый язык. Том V уже построил
спектральный цикл
\begin{equation}
 \mathcal H=\ell^2(\mathbb Z)\otimes q_0\mathbb C^{105},
 \qquad
 Ne_n=ne_n,
 \qquad
 Ue_n=e_{n+1},
 \qquad
 [N,U]=U.
\end{equation}
Спектральная проекция
\begin{equation}
 P=\chi_{[0,\infty)}(N)
\end{equation}
выводит пространство Харди и односторонний сдвиг $S=PUP$. Его компактный
дефект равен
\begin{equation}
 Q_{\rm def}
 =(1-SS^*)\otimes q_0
 =|e_0\rangle\langle e_0|\otimes q_0,
 \qquad
 \rank Q_{\rm def}=15,
 \qquad
 \tau(Q_{\rm def})=\frac17.
\end{equation}
Обратная ориентация даёт сопряжённую пару индексов $-15/+15$.

Это не пространственная трубка. Это конечранговое незамыкание оператора
перехода. Оно имеет ненулевую алгебраическую опору без предположения о
гладком радиусе.

\section{Связь с локальной динамикой}

Минимальный локальный квантовый перенос проекта имеет унитарный символ
\begin{equation}
 W(k)=
 \begin{pmatrix}
 n e^{ik}&-im\\
 -im&n e^{-ik}
 \end{pmatrix},
 \qquad m^2+n^2=1,
\end{equation}
и непрерывный предел
\begin{equation}
 E^2=p^2+M^2.
\end{equation}
При $m=0$ остаются два встречных безмассовых хиральных сдвига с общим
световым конусом. Ненулевая щель $M$ кинематически допустима, но её
амплитуда текущим родителем не выбрана. Дираковские пределы обратимых
квантовых блужданий являются стандартным мостом между дискретным правилом
шага и континуальной волной
\cite{transition-bisio-dariano-tosini2012,
transition-arrighi-forets-nesme2013}.

\section{Новый словарь без новой сущности}

После решения ветви ранние образы читаются так:
\begin{center}
\begin{tabular}{ll}
\toprule
Ранний образ & Спектральное чтение\\
\midrule
нить & упорядоченная композиция локальных переходов\\
виток & степень $U^n$, измеряемая оператором $N$\\
ткань & однородная фаза сети переходов или квантового блуждания\\
узел & локализованный конечранговый дефект оператора перехода\\
частица/античастица & Real-сопряжённая пара ориентаций $-15/+15$\\
масса & щель квазиэнергии локального переноса\\
размер & длина локализации, а не механическая толщина\\
\bottomrule
\end{tabular}
\end{center}

Этот словарь совпадает с уже оформленной первичностью перехода:
соответствия Мориты типизируют стрелки, внутреннее тензорное произведение
составляет пути, а индекс и голономия классифицируют дефекты. Строгий язык
следует теории сильной эквивалентности и расширений операторных алгебр
\cite{transition-brown-green-rieffel1977,
transition-pimsner-voiculescu1980}.

\section{Почему ранг пятнадцать ещё не является атомом}

Проектор ранга пятнадцать в полной комплексной матричной алгебре
алгебраически раскладывается:
\begin{equation}
 q_0=q_1+\cdots+q_{15},
 \qquad \rank q_j=1,
\end{equation}
а его комплексный $K_0$-класс равен пятнадцати генераторам класса один.
Поэтому из одного равенства $\rank Q_{\rm def}=15$ нельзя заключить, что
получена элементарная частица или неделимый квант пространства.

Возможная физическая неделимость должна следовать из более строгого
условия: ни один собственный подпроектор не должен одновременно
сохраняться полной алгеброй, Real-структурой, градуировкой, gauge-действием
и локальным переносом.

\begin{theorem}[Архитектурное решение]
В текущей версии фундаментальной рабочей сущностью выбирается локальный
спектральный переход и его конечранговый дефект, а не непрерывная
пространственная нить. Класс пятнадцать является точным кандидатом
материального пакета, но его физическая примитивность ещё не доказана.
\end{theorem}

\section{Следующий различающий тест}

Следующий гейт
\texttt{version6\_spectral\_transition\_minimal\_support\_gate}
должен классифицировать все подпроекторы $r\le q_0$ и проверить их
инвариантность относительно:
\begin{enumerate}
 \item левого и правого действий конечной алгебры;
 \item Real-инволюции и градуировки;
 \item остаточной gauge- и семейной симметрии;
 \item оператора переноса и тёплицева boundary-класса.
\end{enumerate}

Нулевая гипотеза: полный физический коммутант оставляет $q_0$ неприводимым
и запрещает изолированные ранги $1,3,5$ без нарушения хотя бы одного
условия.

Стоп-критерий: если допустим инвариантный подпроектор меньшего ранга,
класс пятнадцать является составным пакетом, а элементарный спектральный
квант следует искать в его минимальных орбитах.

\begin{protocol}
\begin{enumerate}
 \item скрытая иерархия длины для верёвки: \textbf{не найдена};
 \item буквальная непрерывная нить как активная ветвь: \textbf{закрыта};
 \item спектральный оператор шага $U$ и счётчик $N$: \textbf{существуют};
 \item конечранговый дефект: \textbf{ранг $15$, вес $1/7$};
 \item общий световой конус локального переноса: \textbf{получен};
 \item ненулевая массовая щель: \textbf{не выведена};
 \item алгебраическая примитивность ранга $15$: \textbf{нет};
 \item физическая эквивариантная примитивность: \textbf{открыта};
 \item выбранная парадигма: \textbf{первичность спектрального перехода};
 \item следующая проверка: \textbf{минимальная инвариантная опора}.
\end{enumerate}
\end{protocol}

Машинный сертификат сохранён в
\path{s2t_v6_single_thread_scale_hierarchy_branch_decision_gate_results.json}.

\begin{thebibliography}{9}
\bibitem{transition-bisio-dariano-tosini2012}
A.~Bisio, G.~M.~D'Ariano, A.~Tosini,
\emph{Quantum Field as a Quantum Cellular Automaton: the Dirac Free
Evolution in One Dimension}, arXiv:1212.2839.

\bibitem{transition-arrighi-forets-nesme2013}
P.~Arrighi, M.~Forets, V.~Nesme,
\emph{The Dirac Equation as a Quantum Walk: Higher Dimensions,
Observational Convergence}, arXiv:1307.3524.

\bibitem{transition-brown-green-rieffel1977}
L.~G.~Brown, P.~Green, M.~A.~Rieffel,
\emph{Stable Isomorphism and Strong Morita Equivalence of
$C^*$-Algebras}, Pacific Journal of Mathematics 71 (1977), 349--363.

\bibitem{transition-pimsner-voiculescu1980}
M.~Pimsner, D.~Voiculescu,
\emph{Exact Sequences for $K$-Groups and Ext-Groups of Certain
Cross-Product $C^*$-Algebras}, Journal of Operator Theory 4 (1980),
93--118.
\end{thebibliography}